% Part: many-valued-logics
% Chapter: sequent-calculus
% Section: proof-theoretic-notions

\documentclass[../../../include/open-logic-section]{subfiles}

\begin{document}

\olfileid[pt]{mvl}{seq}{prf}

\olsection{Noções da teoria da prova}

\begin{explain}
Assim como definimos várias noções semânticas importantes (validade,
consequência, satisfatibilidade), agora definimos as correspondentes
\emph{noções da teoria da prova}. Elas não são definidas recorrendo à
satisfação de !!{sentence}s em !!{structure}s, mas recorrendo à
!!{derivability} ou à !!{nonderivability} de certos sequentes. Foi uma
descoberta importante que essas noções coincidem. O fato de
coincidirem é o conteúdo da \emph{correção} e do \emph{teorema da
completude}.
\end{explain}

\begin{defn}[Teoremas]
!!^a{sentence}~$!A$ é um \emph{teorema} se existe !!a{derivation}
em~$\Log{LK}$ do sequente $\quad \Sequent !A$. Escrevemos $\Proves
!A$ se $!A$ é um teorema e $\Proves/ !A$ se não é.
\end{defn}

\begin{defn}[!!^{derivability}]
!!^a{sentence}~$!A$ é \emph{!!{derivable} a partir de} um conjunto de
!!{sentence}s~$\Gamma$, $\Gamma \Proves !A$, sse existe um
subconjunto finito~$\Gamma_0 \subseteq \Gamma$ e uma sequência
$\Gamma_0'$ das !!{sentence}s em~$\Gamma_0$ tais que $\Log{LK}$
!!{derive}s $\Gamma_0' \Sequent !A$. Se $!A$ não é !!{derivable} a
partir de $\Gamma$, escrevemos $\Gamma \Proves/ !A$.
\end{defn}

Por causa das regras de contração, enfraquecimento e troca, a ordem e
o número de !!{sentence}s em~$\Gamma_0'$ não importam: se um sequente
$\Gamma_0' \Sequent !A$ é !!{derivable}, então $\Gamma_0'' \Sequent
!A$ também é, para qualquer $\Gamma_0''$ que contenha as mesmas
!!{sentence}s que~$\Gamma_0'$. Por exemplo, se $\Gamma_0 = \{!B, !C\}$,
então tanto $\Gamma_0' = \tuple{!B, !B, !C}$ quanto $\Gamma_0'' =
\tuple{!C, !C, !B}$ são sequências que contêm apenas as !!{sentence}s
em~$\Gamma_0$. Se um sequente que contém uma delas é !!{derivable}, o
que contém a outra também é, por exemplo:
\begin{prooftree}
  \AxiomC{}
  \Deduce$!B, !B, !C \fCenter !A$
  \RightLabel{\LeftR{\Contraction}}
  \UnaryInf$!B, !C \fCenter !A$
  \RightLabel{\LeftR{\Exchange}}
  \UnaryInf$!C, !B \fCenter !A$
  \RightLabel{\LeftR{\Weakening}}
  \UnaryInf$!C, !C, !B \fCenter !A$
\end{prooftree}
De agora em diante, diremos que, se $\Gamma_0$ é um conjunto finito de
!!{sentence}s, então $\Gamma_0 \Sequent !A$ é qualquer sequente cujo
antecedente seja uma sequência de !!{sentence}s em~$\Gamma_0$, e
incluiremos tacitamente contrações, trocas e enfraquecimentos quando
necessário.

\begin{defn}[Consistência]
Um conjunto de sentenças~$\Gamma$ é \emph{inconsistente} sse existe um
subconjunto finito~$\Gamma_0 \subseteq \Gamma$ tal que $\Log{LK}$
!!{derive}s $\Gamma_0 \Sequent \quad$. Se $\Gamma$ não é inconsistente,
isto é, se, para todo $\Gamma_0 \subseteq \Gamma$ finito, $\Log{LK}$
não !!{derive} $\Gamma_0 \Sequent \quad$, dizemos que é
\emph{consistente}.
\end{defn}

\begin{prop}[Reflexividade]
\ollabel{prop:reflexivity}
Se $!A \in \Gamma$, então $\Gamma \Proves !A$.
\end{prop}

\begin{proof}
O sequente inicial $!A \Sequent !A$ é !!{derivable}, e $\{!A\}
\subseteq \Gamma$.
\end{proof}
  
\begin{prop}[Monotonicidade]
\ollabel{prop:monotonicity}
Se $\Gamma \subseteq \Delta$ e $\Gamma \Proves !A$, então $\Delta
\Proves !A$.
\end{prop}

\begin{proof}
Suponha $\Gamma \Proves !A$, isto é, existe um conjunto finito
$\Gamma_0 \subseteq \Gamma$ tal que $\Gamma_0 \Sequent !A$ é
!!{derivable}. Como $\Gamma \subseteq \Delta$, então $\Gamma_0$ também
é um subconjunto finito de~$\Delta$. A !!{derivation} de $\Gamma_0
\Sequent !A$ também mostra, portanto, que $\Delta \Proves !A$.
\end{proof}

\begin{prop}[Transitividade]
\ollabel{prop:transitivity}
Se $\Gamma \Proves !A$ e $\{!A\} \cup \Delta \Proves !B$, então
$\Gamma \cup \Delta \Proves !B$.
\end{prop}

\begin{proof}
Se $\Gamma \Proves !A$, existe um $\Gamma_0 \subseteq \Gamma$ finito
e !!a{derivation}~$\pi_0$ de $\Gamma_0 \Sequent !A$. Se $\{!A\} \cup
\Delta \Proves !B$, então, para algum subconjunto finito $\Delta_0
\subseteq \Delta$, existe !!a{derivation}~$\pi_1$ de $!A, \Delta_0
\Sequent !B$. Considere a seguinte !!{derivation}:
\begin{prooftree}
  \AxiomC{}
  \RightLabel{$\pi_0$}
  \Deduce$\Gamma_0 \fCenter !A$
  \AxiomC{}
  \RightLabel{$\pi_1$}
  \Deduce$!A, \Delta_0 \fCenter !B$
  \RightLabel{\Cut}
  \BinaryInf$\Gamma_0, \Delta_0 \fCenter !B$
\end{prooftree}
Como $\Gamma_0 \cup \Delta_0 \subseteq \Gamma \cup \Delta$, isso
mostra que $\Gamma \cup \Delta \Proves !B$.
\end{proof}

Observe que isso significa, em particular, que, se $\Gamma \Proves
!A$ e $!A \Proves !B$, então $\Gamma \Proves !B$. Também se segue
que, se $!A_1, \dots, !A_n \Proves !B$ e $\Gamma \Proves !A_i$ para
cada~$i$, então $\Gamma \Proves !B$.

\begin{prop}
\ollabel{prop:incons}
$\Gamma$ é inconsistente sse $\Gamma \Proves {!A}$ para toda
  sentença~$!A$.
\end{prop}

\begin{proof}
Exercício.
\end{proof}

\begin{prob}
Demonstre \olref[fol][seq][ptn]{prop:incons}
\end{prob}

\begin{prop}[Compacidade]
  \ollabel{prop:proves-compact}
  \begin{enumerate}
  \item Se $\Gamma \Proves !A$, então existe um subconjunto finito
    $\Gamma_0 \subseteq \Gamma$ tal que $\Gamma_0 \Proves !A$.
  \item Se todo subconjunto finito de~$\Gamma$ é consistente, então
    $\Gamma$ é consistente.
  \end{enumerate}
\end{prop}

\begin{proof}
  \begin{enumerate}
    \item Se $\Gamma \Proves !A$, então existe um subconjunto finito
      $\Gamma_0 \subseteq \Gamma$ tal que o sequente $\Gamma_0
      \Sequent !A$ tem !!a{derivation}. Consequentemente, $\Gamma_0
      \Proves !A$.
    \item Se $\Gamma$ é inconsistente, existe um subconjunto
      finito~$\Gamma_0 \subseteq \Gamma$ tal que $\Log{LK}$
      !!{derive}s $\Gamma_0 \Sequent \quad$. Mas então $\Gamma_0$ é um
      subconjunto finito inconsistente de~$\Gamma$.
  \end{enumerate}
\end{proof}

\end{document}
